% note revision: acf934d71fce0e4805615b9b9aceb4ea3bdae444
% Generated by make_note_labels.py --rev acf934d71fce; do not edit by hand. The note's labels cited by paper_body.tex, supplemental_body.tex, main_paper.tex and main_supplement.tex, as \newlabel lines of that revision's main_note.aux, prefixed N- as xr-hyper prefixed them.
\makeatletter
\newlabel{N-app:release}{{G}{131}{Result package and reproduction limits}{appendix.G}{}}
\newlabel{N-app:release-edges}{{G.3}{133}{Bin edges}{subsection.G.3}{}}
\newlabel{N-app:response-mismatch}{{D}{128}{Detector-response-mismatch closure test}{appendix.D}{}}
\newlabel{N-sec:205bins}{{A.3.1}{110}{The 205-bin selection}{subsubsection.A.3.1}{}}
\newlabel{N-sec:3d}{{7}{44}{Higher-dimensional extensions: available energy and momentum transfer}{section.7}{}}
\newlabel{N-sec:3d-anchor}{{7.3}{47}{Validation against the 2D measurement}{subsection.7.3}{}}
\newlabel{N-sec:3d-classifier}{{7.5}{48}{Classifier-family cross-check (GBDT vs.\ NN)}{subsection.7.5}{}}
\newlabel{N-sec:controlplots}{{2.5}{11}{Reco-level control distributions}{subsection.2.5}{}}
\newlabel{N-sec:coverage}{{5.2}{32}{Fixed-truth coverage of the 2D statistical band}{subsection.5.2}{}}
\newlabel{N-sec:coverage-5d}{{9.5}{68}{Five-dimensional coverage test}{subsection.9.5}{}}
\newlabel{N-sec:coverage-tunes}{{5.3}{33}{Held-out-tune coverage of the 2D covariance}{subsection.5.3}{}}
\newlabel{N-sec:eavailw-adopted}{{9.6}{69}{Status}{section*.47}{}}
\newlabel{N-sec:eavailw-axis}{{8.1}{55}{$W$ as a fifth axis}{subsection.8.1}{}}
\newlabel{N-sec:eavailw-construction}{{9.6}{69}{The five-axis covariance of the configuration-(a) central value and its $(\Eavail ,W)$ projection}{subsection.9.6}{}}
\newlabel{N-sec:eavailw-generators}{{8.3}{56}{Generator predictions and the target difference}{subsection.8.3}{}}
\newlabel{N-sec:fived16}{{10}{71}{Illustrative sixteen-cell five-observable output on data}{section.10}{}}
\newlabel{N-sec:fived16-fd}{{10.5}{77}{Fake-data tests of the estimator's bias}{subsection.10.5}{}}
\newlabel{N-sec:fived16-gen}{{10.7}{81}{Comparison with generator predictions}{subsection.10.7}{}}
\newlabel{N-sec:fived16-gen-target}{{10.8}{82}{The predictions' target and definitions}{subsection.10.8}{}}
\newlabel{N-sec:fived16-ibu}{{10.6}{80}{Dependence on the unfolding algorithm}{subsection.10.6}{}}
\newlabel{N-sec:fived16-pl}{{10.4}{76}{The run-period block}{subsection.10.4}{}}
\newlabel{N-sec:gbdt-diagnostic-definitions}{{B.2}{119}{Scalar-5D diagnostic reductions and interval definitions}{subsection.B.2}{}}
\newlabel{N-sec:gbdt-model-dependence}{{9.4}{63}{Model dependence, variability and reporting resolution}{subsection.9.4}{}}
\newlabel{N-sec:joint5d}{{11}{88}{Calibrated joint tests of generator predictions}{section.11}{}}
\newlabel{N-sec:joint5d-calibration}{{11.2}{89}{Calibration within the implemented simulation model}{subsection.11.2}{}}
\newlabel{N-sec:joint5d-ch}{{11.9}{97}{Hydrocarbon target correction for NuWro and GiBUU}{subsection.11.9}{}}
\newlabel{N-sec:joint5d-conditions}{{11.4}{91}{Conditions}{subsection.11.4}{}}
\newlabel{N-sec:joint5d-design}{{11.1}{88}{Hypotheses, grid and statistics}{subsection.11.1}{}}
\newlabel{N-sec:joint5d-results}{{11.5}{93}{Decisions}{subsection.11.5}{}}
\newlabel{N-sec:method}{{3}{13}{The unbinned OmniFold procedure}{section.3}{}}
\newlabel{N-sec:method-config}{{3.2}{15}{Estimators and configuration}{subsection.3.2}{}}
\newlabel{N-sec:method-driver}{{3.3}{18}{MINERvA driver and pipeline conventions}{subsection.3.3}{}}
\newlabel{N-sec:negweight-2d}{{C.4}{124}{Comparison on the 2D measurement}{subsection.C.4}{}}
\newlabel{N-sec:negweight-binned}{{C.2}{123}{Binned limit: exact reduction to the purity correction}{subsection.C.2}{}}
\newlabel{N-sec:negweight-footing}{{C.6}{126}{Which background treatment each production chain uses}{subsection.C.6}{}}
\newlabel{N-sec:negweight-uq}{{C.5}{125}{Propagation through systematics and statistics}{subsection.C.5}{}}
\newlabel{N-sec:results}{{6}{35}{Results}{section.6}{}}
\newlabel{N-sec:results-tension}{{6.4}{41}{Eigenmode analysis of the published-covariance tension, $\chi ^{2}=\chiPaper $}{subsection.6.4}{}}
\newlabel{N-sec:systematics}{{4}{20}{Uncertainty quantification: constructing the covariance blocks}{section.4}{}}
\newlabel{N-sec:tension}{{A.5}{114}{Anatomy of the paper-covariance tension}{subsection.A.5}{}}
\newlabel{N-sec:universes}{{A.2.4}{106}{Universes}{subsubsection.A.2.4}{}}
\newlabel{N-sec:uq-2d-scope}{{4.6}{25}{Scope, conventions and limitations of the 2D covariance}{subsection.4.6}{}}
\newlabel{N-sec:uq-combine}{{4.5}{24}{Combination and the unified-throw test}{subsection.4.5}{}}
\newlabel{N-sec:val-departure}{{9.3}{61}{Why the estimator does not follow a generator-sized change}{subsection.9.3}{}}
\newlabel{N-sec:val-negweight}{{9.2}{60}{Background subtraction and generator-sized truth changes}{subsection.9.2}{}}
\newlabel{N-sec:validation}{{5}{28}{Validation of the two-dimensional measurement}{section.5}{}}
\newlabel{N-tab:fived16-fd}{{19}{78}{Fake-data tests on the sixteen cells: $\chi ^2$ (16 degrees of freedom), $p$ and the cells with $|z|>\fvval {trig-z}$ of the bias of the mean against the reported $C_{\mathrm {total}}+C_{\mathrm {resp}}+C_{\mathrm {PL,5D}}$, against $C_{\mathrm {total}}+C_{\mathrm {resp}}$ (without the run-period block), and against $C_{\mathrm {total}}$ alone, each with the truth's own generator shift removed where the second column says so ($-$). $n$ is the number of pseudo-experiments}{table.19}{}}
\newlabel{N-tab:pl-added}{{4}{25}{The run-period block added to the other blocks}{table.4}{}}
\makeatother
